%%%%%%%%%%  scan idx 174  |  folio 161  %%%%%%%%%%
plexes acycliques, est pleinement fidèle et a pour image essentielle $\mathrm D^{-}(\mathcal C_S)_{\mathcal C_0\text{-coh}}$. Si $\mathcal C_{0S}$ est relevable dans $\mathcal C_S$ (I 1.2) (i.e.\ si les objets de $\mathcal C_{0S}$ sont projectifs), la catégorie $\mathrm K^{-,\varnothing}(\mathcal C_{0S})$ est réduite aux objets nuls, et par suite le foncteur
\[
\mathrm K^{-}(\mathcal C_{0S})\longrightarrow \mathrm D(\mathcal C_S)
\]
est pleinement fidèle (et a pour image essentielle $\mathrm D^{-}(\mathcal C_S)_{\mathcal C_0\text{-coh}}$).
\end{enumerate}

\emph{Preuve.} a) Soit
\[
0\longrightarrow L\longrightarrow L^0\longrightarrow\ldots\longrightarrow L^m\longrightarrow0
\]
une suite exacte de $\mathcal C_S$, avec $L^i\in\Ob\mathcal C_{0S}$ pour tout $i$ ; il s'agit de montrer que $L$ est de $\mathcal C_{0S}$-type fini. Or, d'après (I 2.2), $L$ est $\mathcal C_0$-pseudo-cohérent, donc (I 2.9) de $\mathcal C_0$-présentation finie, et l'assertion résulte de (ii).

b) Montrons, par récurrence descendante sur $i$, que $E$ est $i$-$\mathcal C_{0S}$-pseudo-cohérent pour $i\ge n$ (resp. $i\ge n+1$). On peut déjà supposer $E$ strictement $(n+1)$-$\mathcal C_{0S}$-pseudo-cohérent (resp. $(n+2)$-$\mathcal C_{0S}$-pseudo-cohérent) (I 2.1), et il s'agit de prouver que $E$ est $n$-$\mathcal C_{0S}$-pseudo-cohérent (resp. $(n+1)$-$\mathcal C_{0S}$-pseudo-cohérent). On a donc un suite exacte
\[
0\longrightarrow E'\longrightarrow E\longrightarrow E''\longrightarrow0,
\]
où $E'$ est strictement $\mathcal C_{0S}$-parfait (I 2.1) et $E''{}^i=0$ pour $i>n$ (resp. $i>n+1$). Quitte à remplacer $E$ par $E''$, on peut donc supposer (I 2.6) que $E^i=0$ pour $i>n$ (resp. $i>n+1$). Alors, d'après (I 2.10 a)), $H^n(E)$ (resp. $H^{n+1}(E)$) est de $\mathcal C_0$-type fini (resp. de $\mathcal C_0$-présentation finie). Donc, d'après (ii), $H^n(E)$ (resp. $H^{n+1}(E)$) est de $\mathcal C_{0S}$-type fini. Appliquant à nouveau (I 2.10 b)), on en déduit que $E$ est $n$-$\mathcal C_{0S}$-pseudo-cohérent (resp. $(n+1)$-$\mathcal C_{0S}$-pseudo-cohérent), cqfd. (N.B. l'hypothèse (i) et l'assertion a) ont été utilisées implicitement dans les références à (I 2.6 et 2.10)).

c) La première assertion découle de b), la seconde de (I 2.7) ; la dernière est évidente.

\textbf{Proposition 1.2.} Supposons $\mathcal C_0$ localement relevable dans $\mathcal C$ (I 1.2) et localement stable par noyau d'épimorphisme (I 1.2), et faisons les 
hypo%
